%% DRAFT, NOT SUBMITTED. RNAAS draft, AASTeX v7.0.1 (compile on Overleaf; the AAS template ships aastex701.cls).
%% Every number traces to a file in this repository; see salinas-second-transits/rnaas/numbers.md.
%% Hold until (1) the Salinas team has replied or had a fair chance to (outreach/), (2) the blind
%% session's new-sector check is in results/round3/independent/RESULTS.md, (3) the 2026-12-25
%% transit of TIC 107113048 has been looked for, and (4) the AAS AI policy has been re-checked.
\documentclass{aastex701}

\begin{document}

\title{Periods for Four TESS Single-transit Candidates from Later Sectors}

\author[orcid=0009-0000-7524-4470]{Jon McBee}
\affiliation{Independent researcher, USA}
\email{jon.mcbee@gmail.com}

\begin{abstract}
\citet{Salinas2026} listed 108 single-transit candidates and 4 multi-planet systems from TESS
Sectors 27--60. I searched Sectors 61--107 for repeats of each catalogued event. On 112 stars, 53
candidates show at least one matching event and 8 have three or more events on one period. Two
are already TOIs, two are unreliable, and four are new. Their periods are
$17.06402$, $21.19694$, $49.70152$ and $23.75297$\,d. All four pass odd/even, difference-image and
Gaia checks, and a blind re-analysis with separate code recovers each period
within $1.6\sigma$. TIC 241550976 and 107113048 look like warm giants on circular orbits. The
transit of TIC 275647095 is too long for a circular orbit, and a marginal dip at phase 0.87 may make it
an eclipsing binary. TIC 230133616
was never a single: it transits 37 times in the northern continuous viewing zone, and its host is a
subgiant with an unresolved companion.
\end{abstract}

\keywords{\uat{Exoplanet detection methods}{489} --- \uat{Transit photometry}{1709}}

\section{Sample and Search}

\citet{Salinas2026} found 169 candidates in TESS Sectors 27--60 and fitted their transits with a
CNN that gives calibrated uncertainties. I took the 108 single-transit candidates and the 8
single events in 4 multi-planet systems: 116 events on 112 stars. Of these stars, 101 have
light curves from Sector 61 or later, through Sector 107. I used TESS-SPOC \citep{Caldwell2020},
QLP \citep{Huang2020} and SPOC light curves. I also made my own TESScut aperture photometry
\citep{Brasseur2019} for mid-times in Sectors $\geq$94, where QLP times are offset by up to
7\,min.

The search detrends each sector with a biweight filter \citep{Hippke2019}. It then slides a box of
the catalogued duration across the sector. An event passes if it has SNR $\geq7$ (with red noise),
a depth and duration within 0.5--2 times the catalogued values, and an SNR at least 1.5 times
that of the strongest brightening in the sector. I masked four windows in Sectors 94--97 where
most stars dip together. Then 191 events pass, on 53 candidates. Of the 36 candidates with two
or more distinct events, 8 fit one period that includes the catalogued event.

Two of the eight were already TOIs with the same period: TIC 317132251 (TOI-7717.01, 25.21\,d) and
TIC 454039549 (TOI-7743.01, 10.97\,d). Two are unreliable. TIC 141423461 has no clean event, and
its new events are twice the catalogued duration. For TIC 198009543, a 9.7\,h transit on an
11.71\,d orbit would need a star larger than the TIC gives: the shortest circular orbit allowed is
42.8\,d. The other four are not TOIs or CTOIs, have no ExoFOP follow-up, and appear in no paper I
could find.

\section{Vetting}

For each of the four stars I fitted supersampled trapezoids to every clean transit. I took
mid-times from TESS-SPOC, QLP before Sector 94, and my TESScut photometry after that, and depths
from QLP and TESS-SPOC. A weighted linear ephemeris fits the times with reduced $\chi^2$ of
0.3--3.2 and no point beyond $3\sigma$. Every shorter alias predicts transits on flat data, and the
blind analysis below excludes each of them at $\geq15\sigma$.

Odd and even depths agree within $1.3\sigma$ ($2.2\sigma$ in the blind analysis, for the star with
2+2 transits). A box scan at all phases finds no secondary eclipse
once single bad orbits are dropped, except a marginal dip for TIC 275647095 (Section~3). The $3\sigma$ limits at phase 0.5 are 44--705\,ppm.
None of the stars is in the TESS EB catalogue \citep{Prsa2022} or VSX \citep{Watson2006} within
$30''$.

TESScut difference images put each dip 1--4$''$ from the target, 0.1--0.2 pixel and within the
systematic error of the method. Every resolved Gaia DR3 \citep{Gaia2023} neighbour bright enough to
cause the dip is excluded as its source. I compared transit durations with the
stellar density from the Gaia parallax \citep{Seager2003}.

A separate analysis then re-derived every period from the light curves alone. It used new code
written without access to the first analysis. It recovered all four periods within $1.6\sigma$,
and it raised the same concerns independently.

\section{Results}

\textit{TIC 241550976} is the cleanest case. The Gaia radius ($0.96\,R_\odot$) gives a
stellar density that fits a circular orbit. The $\approx1.2\,R_{\rm J}$ companion is most likely a
warm giant planet.

\textit{TIC 107113048} also fits a circular orbit ($\approx0.9\,R_{\rm J}$, $R_\star=1.0\,R_\odot$),
from four transits. The Sector 90 transit falls within 0.7\,h of an eclipse of TIC 107116357.
That star is $97''$ away, and its 7--18\% eclipses are not in Gaia DR3, the TESS EB catalogue or
VSX. The transit is on time to 1\,min, and the binary's other eclipses do not dim the target's
pixels. Still, it is the only transit that rules out 74.55 and 149.10\,d. The next transit, 2026
December 25 21:44\,UTC ($\pm11$\,min), falls outside TESS coverage; Sector 114 (2027 February--March)
can rule out 74.55 and 29.82\,d independently.

\textit{TIC 275647095} has a transit 1.42 times longer than a circular orbit allows for its star
($R_\star=1.14\pm0.06\,R_\odot$ from Gaia). The contamination from a $G=14.3$ star $20''$ away
does not explain it: that star lies 42$\sigma$ from the dip's position. An orbit with
$e\gtrsim0.34$ would explain it. But the blind analysis also finds a dip of $694\pm147$\,ppm at
phase 0.872, in 8 of 9 orbits, with no background or centroid change. The transiting object cannot
make it. If real, it is the secondary eclipse of an eccentric eclipsing binary blended within
$\approx2''$, and the signal is a false positive. Sectors 111--114 cover six of these dips.

\textit{TIC 230133616} lies at ecliptic latitude $+83^\circ$. Its transits appear in Sectors 16--86,
37 in all, so it was never a single-transit candidate. The host is a $2.45\,R_\odot$ subgiant with
Gaia RUWE 3.8, which points to an unresolved companion. The transiting object is
$1.7$--$1.9\,R_{\rm J}$, and a circular orbit around this star fits worse by $\Delta\chi^2\approx49$.
A search over the orbits allowed by the transit shape finds no secondary eclipse. That excludes a
3100\,K M dwarf at $>2\sigma$ for 218 of 228 orbits, but not a blend in the unresolved pair.

\begin{deluxetable*}{lllllll}
\tablecaption{The four new periodic candidates.\label{tab:four}}
\tablehead{\colhead{TIC} & \colhead{$P$ (d)} & \colhead{$T_0$ (BTJD)} & \colhead{Depth (ppm)} &
\colhead{$T_{14}$ (h)} & \colhead{$N_{\rm tr}$} & \colhead{Notes}}
\startdata
241550976 & $21.19694\pm0.00004$ & $2339.4442\pm0.0019$ & 16350 & 5.1 & 5 & $\approx1.2\,R_{\rm J}$; circular fits \\
107113048 & $49.70152\pm0.00017$ & $2312.9429\pm0.0033$ & 8240 & 7.1 & 4 & $\approx0.9\,R_{\rm J}$; EB $97''$ away$^a$ \\
275647095 & $17.06402\pm0.00003$ & $2315.5125\pm0.0018$ & 9130 & 8.4 & 7 & $e\gtrsim0.34$ or EB$^c$ \\
230133616 & $23.75297\pm0.00002$ & $2396.1280\pm0.0007$ & 5880 & 9.3 & 31$^b$ & $1.7$--$1.9\,R_{\rm J}$; RUWE 3.8 \\
\enddata
\tablecomments{$T_0$ is the \citet{Salinas2026} event, refitted. $N_{\rm tr}$ counts the
clean transits fitted. $^a$Its eclipse overlaps the Sector 90 transit. $^b$37 transits in all,
including partial ones. $^c$Marginal dip at phase 0.872; Sectors 111--114 will test it.}
\end{deluxetable*}

\section{Conclusion}

Later TESS sectors give periods for four candidates of \citet{Salinas2026}. Sectors 111--114
will test TIC 275647095 and TIC 107113048 again. Three are
southern stars with $T=12.2$--12.4 and transits 0.8--1.6\% deep, well suited to RV follow-up.
Code and tables are at \url{https://github.com/coracle-us/duo-transits}.

\begin{acknowledgments}
This work uses TESS data from MAST: SPOC \dataset[doi:10.17909/t9-nmc8-f686]{https://doi.org/10.17909/t9-nmc8-f686},
TESS-SPOC \dataset[doi:10.17909/t9-wpz1-8s54]{https://doi.org/10.17909/t9-wpz1-8s54} and QLP
\dataset[doi:10.17909/t9-r086-e880]{https://doi.org/10.17909/t9-r086-e880} light curves and
full-frame images \dataset[doi:10.17909/0cp4-2j79]{https://doi.org/10.17909/0cp4-2j79}; Gaia DR3;
ExoFOP \dataset[doi:10.26134/ExoFOP5]{https://doi.org/10.26134/ExoFOP5}; and the NASA Exoplanet
Archive \dataset[doi:10.26133/NEA13]{https://doi.org/10.26133/NEA13}.
\end{acknowledgments}

\facilities{TESS, Gaia}
\software{lightkurve \citep{Lightkurve2018}, astroquery \citep{Ginsburg2019}, astrocut \citep{Brasseur2019}, wotan \citep{Hippke2019}}

\section*{AI disclosure statement}
I used Anthropic's Claude, an AI assistant, under my supervision to write and run the analysis code
and to draft and edit this note. I reviewed the text, checked every number against the saved
outputs and every reference against ADS and arXiv, and take full responsibility for the content.
The note contains no AI-generated figures.

\begin{thebibliography}{}
\bibitem[Brasseur et al.(2019)]{Brasseur2019} Brasseur, C.~E., Phillip, C., Fleming, S.~W., Mullally, S.~E., \& White, R.~L. 2019, Astrocut: Tools for creating cutouts of TESS images, Astrophysics Source Code Library, ascl:1905.007
\bibitem[Caldwell et al.(2020)]{Caldwell2020} Caldwell, D.~A., Tenenbaum, P., Twicken, J.~D., et al. 2020, RNAAS, 4, 201
\bibitem[Gaia Collaboration et al.(2023)]{Gaia2023} Gaia Collaboration, Vallenari, A., Brown, A.~G.~A., et al. 2023, A\&A, 674, A1
\bibitem[Ginsburg et al.(2019)]{Ginsburg2019} Ginsburg, A., Sip{\H o}cz, B.~M., Brasseur, C.~E., et al. 2019, AJ, 157, 98
\bibitem[Hippke et al.(2019)]{Hippke2019} Hippke, M., David, T.~J., Mulders, G.~D., \& Heller, R. 2019, AJ, 158, 143
\bibitem[Huang et al.(2020)]{Huang2020} Huang, C.~X., Vanderburg, A., P\'al, A., et al. 2020, RNAAS, 4, 204
\bibitem[Lightkurve Collaboration et al.(2018)]{Lightkurve2018} Lightkurve Collaboration, Cardoso, J.~V.~d.~M., Hedges, C., et al. 2018, Lightkurve: Kepler and TESS time series analysis in Python, Astrophysics Source Code Library, ascl:1812.013
\bibitem[Pr{\v s}a et al.(2022)]{Prsa2022} Pr{\v s}a, A., Kochoska, A., Conroy, K.~E., et al. 2022, ApJS, 258, 16
\bibitem[Salinas et al.(2026)]{Salinas2026} Salinas, H., Brahm, R., Olmschenk, G., et al. 2026, arXiv:2610.24882
\bibitem[Seager \& Mall{\'e}n-Ornelas(2003)]{Seager2003} Seager, S., \& Mall{\'e}n-Ornelas, G. 2003, ApJ, 585, 1038
\bibitem[Watson et al.(2006)]{Watson2006} Watson, C.~L., Henden, A.~A., \& Price, A. 2006, SASS, 25, 47
\end{thebibliography}

\end{document}
